\documentclass[12pt]{article}
\usepackage[utf8]{inputenc}
\usepackage[T1]{fontenc}
\usepackage{hyperref}
\title{Suwawa Language Resources for K-12 Curriculum Development}
\author{Compiled by Claude Sonnet 5.5 (claude-sonnet-5-5) for Andr\'as Kornai}
\date{}
\begin{document}
\maketitle

\section{Introduction}
Suwawa (ISO 639-3 swu; Glottocode suwa1241) is an Austronesian language of northern Sulawesi, spoken mainly in Suwawa District, Bone Bolango Regency, Gorontalo Province. Wikipedia gives about 5,000 speakers (2012) and also lists the names Bonda, Bone, Bunda and Toewawa. It belongs to the Gorontalic group of the Gorontalo--Mongondow family. Joshua Project notes a shift among younger people toward Indonesian and Gorontalo Malay. The searches found only overview pages and no teaching materials.

\section{Primary Educational Resources}
No primers, textbooks or school programmes for Suwawa were found in search results.

\section{Linguistic and Reference Materials}
\begin{itemize}
\item Wikipedia, ``Suwawa language'': \url{https://en.wikipedia.org/wiki/Suwawa_language}
\item Wikipedia, ``Bolango language'' (a related language): \url{https://en.wikipedia.org/wiki/Bolango_language}
\item Omniglot, Suwawa page (it says the language is closely related to Mongondow; one source only, not verified): \url{https://mail.omniglot.com/writing/suwawa.htm}
\item Glottolog entry for the Gorontalo group of languages: \url{https://glottolog.org/resource/languoid/id/goro1273.bigmap.html}
\end{itemize}

\section{Religious and Community Resources}
\begin{itemize}
\item Joshua Project, Suwawa in Indonesia: \url{https://joshuaproject.net/people_groups/15142/ID}
\end{itemize}

\section{Digital and Audio Resources}
None found in search results.

\section{Curriculum Development Context}
The language is not known to be taught in schools. Gorontalo-region regional-language education is a possible institutional context, but nothing specific was found. Materials would have to be developed from new documentation work with speakers.

\end{document}